\section{The GerryChain Baseline}
\label{sec:gerrychain-baseline}

\subsection{Comparison to MCMC Disruption}

A natural baseline for ``how much disruption is large?'' is the disruption
from a single \recom{} MCMC step, which randomises the plan by a small
but nonzero amount.
A full \recom{} chain of 10{,}000 steps has $d_{\rm Ham} \approx 1$ from
its start plan in the heuristic scale used here (approximately decorrelated).

We compare the reapportionment-induced boundary disruption to three
baselines:

\begin{table}[h]
\centering
\caption{Boundary disruption comparison across disruption types.
         $d_{\rm Ham}$ is the Hamming distance from the reference plan.}
\label{tab:disruption-comparison}
\begin{tabular}{lcc}
\toprule
Disruption type & $d_{\rm Ham}$ & Relative to reapportion \\
\midrule
Single \recom{} step (TX, $k=38$) & $\approx 0.026$ & $0.11\times$ \\
10 \recom{} steps (TX) & $\approx 0.22$ & $0.96\times$ \\
TX reapportionment ($38\to41$) & 0.23 & $1.0\times$ (baseline) \\
Seed change (worst case, GA) & 0.04 & $0.17\times$ \\
Political redistricting (draft comparison range) & 0.35--0.55 & $1.5$--$2.4\times$ \\
\bottomrule
\end{tabular}
\end{table}

The Texas reapportionment disruption ($d_{\rm Ham} = 0.23$) is equivalent
to approximately 10 one-step \recom{} units under the $1/k$ approximation.
This is larger than the typical seed-change disruption (B.7) but
smaller than the political-redistricting comparison range used in this draft.

\subsection{Interpretation}

The finding that reapportionment disruption $\approx 10$ MCMC steps
has an important implication for statistical continuity:

\begin{enumerate}
\item \textbf{Algorithmic reapportionment is smaller than full MCMC remixing in this scale.}
      A certified \recom{} ensemble (10{,}000+ steps) fully decorrelates
      from its starting plan.
      A reapportionment-driven redistricting changes only $\approx$23\% of
      tracts (for the most disruptive case, TX).
      The previous map is a useful starting point for the new map.
\item \textbf{Political-redistricting comparisons need same-metric evidence.}
      Political redistricting cycles (which change 35--55\% of tract
      assignments in the draft comparison range) would produce
      $1.5$--$2.4\times$ more disruption than the Texas scenario, but this
      claim needs either a citation or an enacted-plan calculation.
\end{enumerate}

\subsection{Incumbent Protection Analysis}

One concern about reapportionment is that it disrupts incumbents by
moving them into different districts.
The Hamming distance measure is not an incumbent-disruption measure.
If $d_{\rm Ham} = 0.23$, then 23\% of tracts change district labels after
canonical matching.  Incumbent disruption additionally depends on incumbent
home locations, candidate choices, and district numbering.  Those facts are
outside the current package.

For comparison to incumbent displacement, a future package should geocode
incumbent residences and evaluate whether each residence remains in the
best-matched successor district.  Tract Hamming distance alone is insufficient.

\subsection{GerryChain Baseline Definition}

The ``GerryChain baseline'' in the title of this section refers to using
the \recom{} step size ($d_{\rm Ham} \approx 1/k$ per step) as the unit
of measurement for boundary disruption.
This baseline is natural because:
\begin{enumerate}
\item It is the smallest unit of plan change in the MCMC framework.
\item It is computable for any state and any $k$.
\item It provides an apples-to-apples comparison between algorithmic
      stability and MCMC mixing.
\end{enumerate}

Under this baseline, the 2030 reapportionment for Texas is equivalent to
approximately $0.23 / 0.026 \approx 9$ \recom{} steps.
This is a concrete, policy-relevant characterisation when kept at the right
level: the Texas scenario is a ``9-step-equivalent'' disruption in this MCMC
scale, pending empirical ReCom-chain calibration.
